\section*{Methods}

\subsection*{Spectral representation for the delay model}

Let \(U(t)\) denote the upstream process, such as latent infection incidence, and let \(D(t)\) denote a downstream surveillance observable, such as reported cases, hospital admissions, or deaths. The delay from upstream infection to downstream observation is described by a nonnegative delay density \(g(\tau)\), supported on \(\tau\ge0\) and normalized so that \(\int_0^\infty g(\tau)\,d\tau=1\). A conversion rate \(p\in(0,1]\) represents the probability that an upstream event generates a downstream observation. One common example is the ascertainment rate, which describes the proportion of latent infections ultimately recorded as confirmed cases. Because testing- or reporting-related changes typically occur on coarse temporal scales \cite{colman2023ascertainment,eales2023dynamics,ward2024real} relative to the support of \(g\), we treat \(p\) as locally constant.

We combine delay and conversion into the observation kernel
\[
k(\tau)=p\,g(\tau).
\]
The downstream conditional mean is
\begin{equation}
m(t)
=
\mathbb E\{D(t)\mid U\}
=
(U*k)(t),
\label{eq:delay-mean}
\end{equation}
where \(*\) denotes convolution,
\[
(U*k)(t)=\int_0^\infty U(t-\tau)k(\tau)\,d\tau .
\]
Equivalently, the downstream observation can be written as
\begin{equation}
D(t)=m(t)+\varepsilon(t),
\qquad
\mathbb E\{\varepsilon(t)\mid U\}=0.
\label{eq:mean-residual}
\end{equation}
The observation model is illustrated in Fig.~\ref{fig:intro}A--C.

For additive Gaussian observation models, \(\varepsilon(t)\) corresponds directly to Gaussian noise. For count-based observation models such as Poisson or negative-binomial sampling, equation~\eqref{eq:mean-residual} should instead be interpreted as a mean-plus-residual decomposition, where \(\varepsilon(t)\) denotes deviations from the conditional mean.

Because convolution in the time domain corresponds to multiplication in the frequency domain, it is natural to examine the conditional mean spectrally. Taking Fourier transforms gives
\begin{equation}
\widehat m(f)
=
\widehat U(f)\widehat k(f)
=
p\,\widehat U(f)\widehat G(f),
\label{eq:fourier-delay}
\end{equation}
where \(\widehat X(f)\) denotes the Fourier transform of \(X(t)\). When \(p=1\), equation~\eqref{eq:fourier-delay} reduces to the delay-only setting without conversion loss. This frequency-domain representation highlights the attenuating effect of delay processes: for many epidemiological delay distributions, the magnitude \(|\widehat G(f)|\) decreases with frequency, thereby attenuating high-frequency upstream variation before it reaches the downstream mean. This attenuation mechanism is illustrated in Fig.~\ref{fig:intro}D--F.

This framework extends naturally to multilayer delay structures in which several delay mechanisms act sequentially, such as incubation, testing, reporting, and hospitalization delays. Let \(g_1,g_2,\ldots,g_L\) denote normalized delay kernels associated with successive stages, and let \(p_1,p_2,\ldots,p_L\) denote the corresponding conversion rates. The overall conversion rate is \(p_{\mathrm{tot}}=\prod_{\ell=1}^L p_\ell\), while the overall delay kernel is
\[
g_{\mathrm{tot}}=g_1*g_2*\cdots*g_L.
\]
The overall observation kernel is
\[
k_{\mathrm{tot}}(\tau)=p_{\mathrm{tot}}\,g_{\mathrm{tot}}(\tau),
\]
and the downstream conditional mean satisfies
\[
\mathbb E\{D(t)\mid U\}=(U*k_{\mathrm{tot}})(t).
\]
In the frequency domain, the corresponding transfer function factorizes as
\begin{equation}
\widehat k_{\mathrm{tot}}(f)
=
\prod_{\ell=1}^L p_\ell
\prod_{\ell=1}^L \widehat G_\ell(f).
\end{equation}
Thus, sequential delay mechanisms compound their spectral attenuation multiplicatively, progressively suppressing high-frequency upstream variation before it reaches the downstream process.

\subsection*{Spectral separability of upstream trajectories}

We quantify upstream identifiability by asking whether two candidate upstream trajectories \(U_0\) and \(U_1\) could be distinguished from the downstream observations. The statistical scale for this comparison is set by the downstream likelihood, but the effect of delay is most transparent after expressing the resulting likelihood separability in spectral coordinates. Let
\[
H_0: U=U_0,
\qquad
H_1: U=U_1
\]
denote the two simple hypotheses, and let
\[
m_i(t)=(U_i*k)(t),
\qquad i\in\{0,1\},
\]
be the corresponding downstream mean trajectories. Their observable contrast is
\[
\Delta m(t)=m_1(t)-m_0(t)=\{(U_1-U_0)*k\}(t).
\]

On a finite observation window, let \(\mathbf F\) be the unitary discrete Fourier transform and let \(\mathcal K\) be the diagonal transfer operator with entries \(\widehat k(f_k)=p\widehat G(f_k)\). In the corresponding finite-window spectral representation,
\[
\Delta\mathbf m
=
\mathbf F^*
\mathcal K
\Delta\widehat{\mathbf U},
\]
where \(\Delta\widehat{\mathbf U}\) is the Fourier transform of \(U_1-U_0\). Let \(\mathcal M\) index the downstream observation model, and let \(\mathbf W_{\mathcal M}\) denote the likelihood weight induced by that observation model for the finite contrast between \(m_0\) and \(m_1\). This weight plays the role of observation precision: downstream contrasts in regions with greater observation variability receive smaller weight, and locally the weight reduces to inverse conditional variance. We define the spectral separability between \(U_0\) and \(U_1\) as the likelihood-weighted downstream contrast expressed in the upstream frequency domain,
\begin{equation}
J_{\mathcal M}(U_0,U_1)
=
\Delta\widehat{\mathbf U}^{*}
\mathbf S_{\mathcal M}(U_0,U_1)
\Delta\widehat{\mathbf U},
\label{eq:J-spectral}
\end{equation}
with
\begin{equation}
\mathbf S_{\mathcal M}(U_0,U_1)
=
\mathcal K^{*}
\mathbf F\mathbf W_{\mathcal M}(U_0,U_1)\mathbf F^{*}
\mathcal K.
\label{eq:S-operator}
\end{equation}
The transfer operator \(\mathcal K\) attenuates upstream variation according to the delay distribution and conversion probability. After transforming back to the observation window, \(\mathbf W_{\mathcal M}\) scores the resulting downstream mean contrast on the scale of observation variability specified by the observation model. Consequently, \(J_{\mathcal M}\) is small when upstream differences are concentrated at temporal scales suppressed by the delay kernel, or when the delay-filtered downstream differences are small on the observation-variability scale.

The same quantity can be written on the downstream mean scale as
\begin{equation}
J_{\mathcal M}(U_0,U_1)
=
\Delta\mathbf m^*
\mathbf W_{\mathcal M}(U_0,U_1)
\Delta\mathbf m .
\label{eq:J-time}
\end{equation}
Thus \(J_{\mathcal M}\) measures the part of the upstream spectral contrast that remains after delay filtering and is informative under the downstream observation model. \(\mathbf S_{\mathcal M}\) is a spectral operator that combines delay attenuation with the likelihood weight induced by the observation model. This form retains the intuition that identifiable upstream contrast must both survive delay filtering and remain large relative to observation variability. In special cases such as stationary Gaussian noise, the operator reduces to a frequency-by-frequency weight.

\subsection*{Observation models and non-identifiability bound}

We evaluated three observation models commonly used for epidemic surveillance data: additive Gaussian noise, Poisson counts, and negative-binomial counts. The corresponding likelihood weights are summarized in Table~\ref{tab:noise-models}; full derivations are given in Appendix~\ref{sec:Appendix_S1}. For Gaussian observations with common covariance \(\mathbf C_\varepsilon\), \(\mathbf W_{\mathrm G}=\mathbf C_\varepsilon^{-1}\). If the Gaussian noise is stationary, the Fourier basis diagonalizes this covariance and equation~\eqref{eq:J-spectral} becomes the familiar frequency-wise weighted contrast. For conditionally independent Poisson and negative-binomial counts, \(\mathbf W_{\mathcal M}\) is diagonal in time, with entries determined by the exact likelihood curvature between \(m_0(t)\) and \(m_1(t)\). The local approximations in Table~\ref{tab:noise-models} give the corresponding inverse-variance weights in the small-contrast limit.

With these likelihood-induced weights, the spectral separability in equation~\eqref{eq:J-spectral} is
\begin{equation}
J_{\mathcal M}(U_0,U_1)
=
2D_{\mathrm{KL}}
\left(
P_0^{\mathcal M}\Vert P_1^{\mathcal M}
\right),
\label{eq:J-KL}
\end{equation}
where \(P_i^{\mathcal M}\) is the distribution of the downstream observation vector implied by \(U_i\) under observation model \(\mathcal M\). This representation connects the spectral separability criterion to a likelihood-based testing bound.

For any decision rule \(\phi\), define the total testing error
\[
R(\phi) = P_0(\phi=H_1)+P_1(\phi=H_0).
\]
For two simple hypotheses,
\[
\inf_\phi R(\phi)=1-\lVert P_0-P_1\rVert_{\mathrm{TV}}.
\]
Here, \(\lVert \cdot\rVert_{\mathrm{TV}}\) denotes the total variation distance. Combining equation~\eqref{eq:J-KL} with Pinsker's inequality gives
\begin{equation}
\inf_\phi R(\phi)
\ge
1-\frac12\sqrt{J_{\mathcal M}(U_0,U_1)}.
\label{eq:testing-bound}
\end{equation}
Consequently, for any target total testing error level \(\alpha\in(0,1)\),
\begin{equation}
J_{\mathcal M}(U_0,U_1)
\le
4(1-\alpha)^2
\quad\Longrightarrow\quad
\inf_\phi R(\phi)\ge \alpha .
\label{eq:nonidentifiability-threshold}
\end{equation}
When this condition holds, no statistical procedure can reliably distinguish the two upstream trajectories from the downstream observations alone at the specified error criterion.

\begin{table*}[!t]
\centering
\caption{ \textbf{Likelihood weights for downstream separability.} The separability can be written as \(J_{\mathcal M}=\Delta\mathbf m^*\mathbf W_{\mathcal M}\Delta\mathbf m\), where \(\Delta\mathbf m\) is the delay-filtered downstream contrast. For count models, \(m_s(t)=m_1(t)+s\{m_0(t)-m_1(t)\}\), \(0\le s\le1\). Local approximations apply when \(m_0(t)\) and \(m_1(t)\) are close to a common mean \(\bar m(t)\). }
\label{tab:noise-models}
\begin{tabular}{llll}
\hline
Observation model
&
Variance structure
&
Likelihood weight
&
Local approximation
\\
\hline

Gaussian
&
\(\mathbf C_\varepsilon\)
&
\(\mathbf W_{\mathrm G}=\mathbf C_\varepsilon^{-1}\)
&
\(1/S_\varepsilon(f_k)\) if stationary
\\[8pt]

Poisson
&
\(V_{\mathrm P}(m)=m\)
&
\(\displaystyle
[\mathbf W_{\mathrm P}]_{tt}
=2\int_0^1\frac{1-s}{m_s(t)}\,ds\)
&
\(\displaystyle 1/\bar m(t)\)
\\[10pt]

Negative binomial
&
\(\displaystyle V_{\mathrm{NB}}(m)=m+\frac{m^2}{\kappa}\)
&
\(\displaystyle
[\mathbf W_{\mathrm{NB}}]_{tt}
=2\int_0^1
\frac{1-s}{m_s(t)+m_s(t)^2/\kappa}\,ds\)
&
\(\displaystyle
1/\{\bar m(t)+\bar m(t)^2/\kappa\}\)
\\[10pt]

\hline
\end{tabular}
\end{table*}
